\section{Внесённые исправления}
На основании анализа исходного решения были внесены исправления:
\begin{itemize}
    \item В \texttt{main.cpp} включён быстрый режим потоков: \texttt{std::ios::sync\_with\_stdio(false)} и \texttt{std::cin.tie(nullptr)}.
    \item Вывод сделан буферизованным: ответы на команды накапливаются и сбрасываются порциями по 1 MiB.
\end{itemize}

\subsection{Консоль (исправленная версия)}
\begin{alltt}
root$ g++ -O2 -std=c++20 -Wall -Wextra -pedantic main.cpp -o app_opt
root$ time ./app_opt < bench_input.txt > /dev/null

real    0m4.017s
user    0m2.134s
sys     0m0.141s
root$ g++ -O2 -pg -std=c++20 -Wall -Wextra -pedantic main.cpp -o app_opt_pg
root$ ./app_opt_pg < bench_input.txt > /dev/null
root$ gprof ./app_opt_pg gmon.out > gprof_opt.txt
root$ valgrind --tool=memcheck --leak-check=full --show-leak-kinds=all --show-reachable=no ./app_opt < bench_input.txt > /dev/null 
==14620== Memcheck, a memory error detector
==14620== Copyright (C) 2002-2024, and GNU GPL'd, by Julian Seward et al.
==14620== Using Valgrind-3.24.0 and LibVEX; rerun with -h for copyright info
==14620== Command: ./app_opt
==14620== 
==14620== 
==14620== HEAP SUMMARY:
==14620==     in use at exit: 122,880 bytes in 6 blocks
==14620==   total heap usage: 1,819,698 allocs, 1,819,692 frees, 55,352,697 bytes allocated
==14620== 
==14620== LEAK SUMMARY:
==14620==    definitely lost: 0 bytes in 0 blocks
==14620==    indirectly lost: 0 bytes in 0 blocks
==14620==      possibly lost: 0 bytes in 0 blocks
==14620==    still reachable: 122,880 bytes in 6 blocks
==14620==         suppressed: 0 bytes in 0 blocks
==14620== Reachable blocks (those to which a pointer was found) are not shown.
==14620== To see them, rerun with: --leak-check=full --show-leak-kinds=all
==14620== 
==14620== For lists of detected and suppressed errors, rerun with: -s
==14620== ERROR SUMMARY: 0 errors from 0 contexts (suppressed: 0 from 0)
\end{alltt}

\section{Анализ исправленной версии}
\subsection{Результаты исправленной версии}
\begin{itemize}
    \item Время выполнения: \textbf{real 4.017 s}, \textbf{user 2.134 s}, \textbf{sys 0.141 s}.
    \item \texttt{gprof}: уменьшилось относительное время на обработку операций, связанных с I/O.
    \item \texttt{memcheck}: утечек памяти и ошибок не выявлено.
\end{itemize}

\subsection{Фрагмент gprof (исправленная версия)}
\begin{alltt}
Flat profile:
  %   cumulative   self              self     total           
 time   seconds   seconds    calls  ms/call  ms/call  name    
 90.76      1.37     1.37                             main
  2.65      1.41     0.04   128024     0.00     0.00  RBTree::insertBalance(RBTree::Node*)
  1.99      1.44     0.03    61266     0.00     0.00  RBTree::removeBalance(RBTree::Node*, Side)
  1.99      1.47     0.03                             _init
  1.32      1.49     0.02       13     1.54     1.54  frame_dummy
  0.66      1.50     0.01   700001     0.00     0.00  operator>>(std::istream&, Command&)
\end{alltt}

\subsection{Сравнение}
На выбранном тесте ускорение составило \(\frac{6.057 - 4.017}{6.057}\approx 33.7\%\) по \texttt{real time}.
\pagebreak

